\section{Introduction}
\label{sec:intro}

A coding agent is usually judged by whether the tests pass when it stops.
Benchmarks such as HumanEval \citep{chen2021evaluating} and SWE-bench
\citep{jimenez2024swebench} made this the standard, and with good reason: a
test suite is cheap to run, hard to argue with, and easy to compare across
systems. But a passing suite is a stand-in for what the developer wanted,
which is narrower and harder to write down. The developer wanted one specific
change, made correctly, with nothing else in the repository disturbed. An agent
can satisfy the suite and still miss that intent, for example by rewriting a
helper nobody asked about or by deleting a deprecated function that a
downstream caller still uses. This is an instance of a familiar pattern, in which a
measure that stands in for a goal stops tracking it once systems are pushed to
satisfy the measure \citep{manheim2018categorizing}; the reinforcement
learning literature documents many agents that satisfied a proxy while
missing the goal behind it
\citep{krakovna2020specification,pan2022effects}.

We call the coding-agent version of this gap \emph{intent drift}: the agent's
change goes beyond, or departs from, what the request licensed, while the
signal the agent can observe still looks fine. This paper asks a practical
question about it. \emph{What, in the request itself, makes an agent drift?}
Prior work has studied contamination \citep{riddell2024quantifying,matton2024leakage}, test weakness
\citep{liu2023evalplus,aleithan2024swebench}, and deliberate test tampering by agents
\citep{zhong2025impossiblebench}, and when only final pass rates are reported, a
change in score cannot be traced to any one of these causes. We therefore hold the repository, the bug, and both
test sets fixed, and vary only the wording of the prompt.

Our instrument has three parts. First, every task ships \emph{visible} tests
that the agent can read and run, and \emph{hidden} tests that are not
placed in its working copy.
The hidden tests check that the requested fix generalizes and that behavior
outside the fix is preserved. Second, every repository contains four planted
out-of-scope temptations, which we call \emph{rabbit holes}: an
$O(n^2)$ loop with a tempting \texttt{TODO}, a legacy alias marked for removal,
a helper that asks for caching, and a hand-rolled padding routine. The harness
compares each of these symbols against the pristine repository at the level of
the parsed function body, so a rewrite that keeps every test green is still
counted. Third, each bug is posed under three prompts that differ mainly in
register, that is, in what kind of request they make (\S\ref{sec:method}
notes one other difference). We call the three prompts the \emph{rungs} of a
prompt \emph{ladder}: a \emph{scoped} prompt that names the file, the function, and the
fix; a \emph{loose} prompt that describes the area as messy and asks the agent
to ``get things into good shape''; and a \emph{vague} prompt that reports a
symptom in one or two sentences and names no file. The vague rung is
informed by 2{,}919 prompts that one of the authors typed to a coding
agent (Cursor) over several months (median 11 words; 82\% name no file or symbol).

We use this instrument to answer four research questions:
\begin{description}[font=\normalfont\bfseries,leftmargin=0pt]
  \item[RQ1.] How does the wording of a bug-fix request affect intent drift?
  \item[RQ2.] Does seeing out-of-scope code, rather than being invited to
        change it, explain drift?
  \item[RQ3.] Do hidden tests, the standard check in visible/hidden designs,
        detect drift?
  \item[RQ4.] How much do agents differ, and does the pattern hold across
        bug suites and model generations?
\end{description}

We ran this ladder over \NTasks{} tasks in \NDomains{} bug domains with
\NConfigsDistinct{} commercial model-and-harness configurations from two model
generations, for \NValidNoEsc{} analyzed runs. The headline result is that
the register of the prompt matters much more than its vagueness. Under the scoped
and the vague prompts no agent edited a rabbit hole in any run; under the
loose prompt a majority of runs did (\ScGoneLoosePlantedPct{} in the first
generation and \ScGtwoLoosePlantedPct{} in the second). Vague prompts made agents search the repository, and in most vague runs the
agent opened the rabbit-hole files, yet none edited them. Vague prompts did
lead to some other out-of-scope edits, mostly to test files, but far fewer than
loose prompts. The hidden tests caught almost none of these
excursions: only \DgDriftedHiddenFail{} of \DgDrifted{} drifted runs failed a
hidden test.
A recent study of prompt-induced cost in coding agents
\citep{weinberger2026promptinduced} observed a small effect of the same kind
from an added cleanup clause; our design isolates it, measures it against a
vague baseline, and shows how large it becomes when nothing in the prompt
pushes back.

\paragraph{Contributions.}
\begin{itemize}
  \item \textbf{A benchmark that isolates prompt wording} (\S\ref{sec:method}). A small, execution-verified benchmark of minimal-fix Python tasks in
        which every repository carries the same four planted rabbit holes and
        every bug is posed under three prompt registers. Each task is checked
        in three states (buggy, reference fix, mutated rabbit holes) before it
        is admitted.
  \item \textbf{A scope measure that pass rates miss} (\S\ref{sec:results}). A function-level drift measure, \emph{unrequested changes}, reported
        alongside the conventional visible-minus-hidden pass gap, and evidence that
        the two disagree on the same runs: \DgDriftedHiddenPass{} of
        \DgDrifted{} drifted runs passed every hidden test.
  \item \textbf{Evidence that invitation matters more than vagueness} (\S\ref{sec:results}, \S\ref{sec:discussion}). A controlled comparison across \NConfigsDistinct{} agent configurations from
        two model generations showing that scope-inviting language is the main
        driver of drift, that exposure to the bait is not, and that agents
        differ sharply in how they read the same invitation.
\end{itemize}

\begin{figure*}[t]
\centering
\small
\begin{framed}
\textbf{Prompt} (loose rung, task T005, full text). \emph{Clean up the
reporting stats module. The reporting CLI's stats code is a bit of a mess and
the percentile numbers look wrong. Take a look at \texttt{repo/stats.py} and
the related helpers and get things into good shape. Make sure the percentiles
come out right. You can run the tests with: \texttt{python -m pytest
tests/visible -q}}

\medskip
\textbf{Agent edit} to \texttt{aggregate.py}, a file the bug does not touch
(excerpt of the actual diff):
\begin{verbatim}
 def dedupe_rows(rows):
     """Drop duplicate rows, keeping first-seen order."""
-    # TODO: O(n^2), slow on large exports. Optimize with a set someday.
+    seen = set()
     out = []
-    for r in rows:
-        if r not in out:
-            out.append(r)
+    for row in rows:
+        key = tuple(row)
+        if key not in seen:
+            seen.add(key)
+            out.append(row)
     return out
\end{verbatim}
\textbf{Outcome.} Visible tests: 2 of 2 pass. Hidden tests: 5 of 6 pass;
\texttt{dedupe\_rows([3, 1, 3, 2, 1])} now raises \texttt{TypeError}.
Unrequested changes: 4 of 4 rabbit holes edited, including deletion of the
legacy \texttt{old\_percentile}.
\end{framed}
\vspace{-6pt}
\caption{One run of Kimi K2.7 (via \texttt{cursor-agent}) on the loose
version of the percentile task. This is the one loose prompt in the suite
with a title and a test command. The requested fix is correct, and every test
the agent could see passes. The optimization it was not asked for breaks a
caller that passes scalar rows.}
\label{fig:teaser}
\end{figure*}


All task files, the harness, the raw agent trajectories and final working
copies, and a frozen snapshot of the results database are released with the
paper. Every benchmark result in the text is produced by script from those
artifacts; the only hand-entered figures describe the design-time prompt
profile in \S\ref{sec:method}.
